% 1. Introduction (1/2 page2): What this project is about, including the goal of the project.

This document aims to outline the steps involved in implementing a client-side connector for the DS-SIM server-side applications. DS-SIM, a dsitributed system simulator, is used to model and study the scheduling of distributed systems. The primary objective was to develop a client component that can communicate effectively with the DS-SIM server, manage job scheduling, and implement efficient algorithms for resource allocation.

The project was structured in phases, starting with establishing a reliable handshake mechanism between the client and the server. This initial step ensured that the client could authenticate and establish a connection with the DS-SIM server. Following this, the focus shifted to implementing a highest core scheduling approach, where the client selects servers based on their core count to handle incoming jobs. The final phase involved adopting a "first capable" strategy, which aims to assign jobs to the first server that meets the job's resource requirements. Each phase was critical in building a robust client-side application capable of handling various job scheduling scenarios.

The goal of this project is to deliver a functional client component that can interact with DS-SIM efficiently. This includes handling an authentication handshake, processing jobs in a loop until a specific exit command is received, and implementing a first capable approach or designing an algorithm that outperforms the first capable method. The client must be resilient, able to handle different types of job requests, and capable of making real-time decisions based on the server's resource availability.

Throughout the semester, the project was developed incrementally, with each feature being implemented, tested, and refined. This iterative development process was supported by regular feedback from our tutor, ensuring that the project stayed on track and met the required specifications. The use of version control systems like GitHub played a crucial role in managing the project's progress. Commits were made weekly, documenting each new feature and bug fix, which helped in tracking development and maintaining a clear project history.

The timeframe for this project spanned several weeks over a semester. This extended period allowed for thorough testing and debugging, ensuring that the final client component was both reliable and efficient. The phased approach to development facilitated a deeper understanding of each aspect of the project, from basic communication protocols to advanced scheduling algorithms.